\documentclass[11pt]{memoireuqamV2ang}
%%%%%%%%%%%%%%%%%%%%%%%%%%%%%%%%%%%
% OVERLEAF: Utiliser TexLive 2021 %
%%%%%%%%%%%%%%%%%%%%%%%%%%%%%%%%%%%

% \usepackage{natbib}  % optionnel, doit être chargé avant babel
\usepackage[utf8]{inputenc} % Pour utiliser les caractères accentués
\usepackage[T1]{fontenc}
\usepackage{url}
\usepackage[a-1b]{pdfx}
\usepackage{hyperref}
\usepackage{lipsum}
\usepackage{layout}
\makeatletter
\renewcommand*{\lay@value}[2]{%
  \strip@pt\dimexpr0.351459\dimexpr\csname#2\endcsname\relax\relax mm%
}
\makeatother

% optionnel
\usepackage{graphicx}% Pour les figures
\usepackage{standalone}
\usepackage{tikz}
\usepackage{pgfplots}
\pgfplotsset{compat=1.18} 

% police similaire à calibri, optionnelle (voir https://guidemt.uqam.ca/gabarit-et-mise-en-page/#ressource-697)
\usepackage[sfdefault,lf]{carlito}

%%%%
% Visitez http://www.labmath.uqam.ca/presentation-de-la-classe-memoireuqamV2/ pour plus details
%%%%

%%%%%%%%%%%%
% Acronymes et symboles
%%%%%%%%%%%%
\usepackage[toc,acronym]{glossaries}
\defglsentryfmt[main]{\textbf{\glsgenentryfmt}}
\newglossary[slg]{symbols}{syi}{syg}{Notation}
\loadglsentries{glossaries/notation}
\loadglsentries{glossaries/acronyms}
\loadglsentries{glossaries/glossary}
\makenoidxglossaries

\begin{document}%\layout


%%%% PAGE TITRE FRANÇAIS %%%%%
% Le titre
\title{Mon titre}
% Votre nom complet tel qu'il apparaît à votre dossier du registrariat de l'UQAM
\author{Mon nom}
% Année et mois courant sauf si spécifié autrement par \degreemonth et \degreeyear
\degreemonthfrancais{mois du dépôt}
\degreeyear{année du dépôt}
\uqammemoire %% ou \uqamthese ou \uqamrapport ou \uqamprop 
\matiere{mathématiques}
\thispagestyle{empty}        % La page titre n'est pas numérotée
\maketitle
%%%%%%%%%%%%%%%%%%

%%%% PAGE TITRE ANGLAIS %%%%%

\titleanglais{My thesis}
\degreemonth{month of submission}
\matiereanglais{mathematics}
\maketitleanglais


%%%%%%%%%%%%%%%%%%%%
% Page préliminaires
%%%%%%%%%%%%%%%%%%%%

\pagenumbering{roman} % numérotation des pages en chiffres romains
\stepcounter{page} % sauter pages titre
\stepcounter{page}

\include{remerciements}
\tableofcontents % Pour générer la table des matières
\listoffigures% Pour générer la liste des figures
\listoftables % Pour générer la liste des tableaux

% Ici, on suppose une liste d'acronymes et une fiche de notation. Supprimez ces lignes si elles sont sans objet.
\glsaddall[types={acronym}]
\printnoidxglossary[type=acronym,toctitle=ACRONYMES,nonumberlist,sort=def]
\glsaddall[types={symbols}]
\printnoidxglossary[type=symbols,style=tree,nonumberlist,sort=def,toctitle=NOTATION]
\glsresetall

% Vous devez inclure un résumé en français ainsi que le résumé en français
% Utilisez l'environnement resume pour rédiger le résumé en français et abstract pour rédiger le résumé en francais
\include{gabarit_memoire_en/resume}
\include{gabarit_memoire_en/abstract}

%%%%%%%%%%%%%%%%%%%%
% Document principal
%%%%%%%%%%%%%%%%%%%%
\pagenumbering{arabic}

% Utilisez l'environnement  introduction pour rédiger votre introduction
\include{introduction}

\include{chapter1}
\include{chapter2}

% Utilisez l'environnement  conclusion pour rédiger votre conclusion
\include{conclusion}

%%%%%%%%%%%%%%%%%%%%
% Page liminaires
%%%%%%%%%%%%%%%%%%%%
\appendix % À partir de cette ligne toutes les commandes \chapter subséquentes seront des appendices
\include{appendix}
% index
\printnoidxglossary[toctitle=INDEX,title=INDEX]

\bibliographystyle{apalike-uqam}
\bibliography{references}
\end{document}
